\section{换序：两个操作能不能交换顺序}
\label{sec:17-Synthesis/02-Recurring-Ideas/05}

一个排队系统的仿真，要报告``最长等待时间''。两个人各写了一版。一版跑一万次仿真，每次取出这一次里最长的那个等待，再对一万个数求平均；另一版先把一万次仿真里同一时刻的等待求平均，得到一条平均曲线，再取这条曲线的最大值。数字差了将近两成。

会上争了半小时谁写错了。两个都没写错，它们算的是两个不同的量：\(\E[\max]\) 和 \(\max \E\)。而这两个量之间有一个方向确定的不等号，\(\E[\max]\ge\max\E\)，恒成立，与仿真质量无关。取最大是一个凸操作，它和期望不可交换，交换之后的偏差只朝一个方向。谁该被报告出去取决于要回答的问题是``典型的一天最坏能坏到什么程度''还是``平均而言最忙的时刻有多忙''，这是业务问题；但两个数不该相等，这是数学问题。

这条缝在这本书里出现了太多次，而且几乎每次都是悄悄跨过去的。

\subsection{一个动作}

第五个反复出现的动作，是这五个里最像记账的一个：\emph{把两个操作交换顺序。}极限与积分、导数与期望、求和与求和、期望与最大、梯度与采样——这些交换在前面十六个部分被使用了无数次，大多数时候连一句说明都没有，因为大多数时候它们确实成立。麻烦在于不成立的时候没有任何东西会报错，你会得到一个类型正确、量纲正确、看上去正常的数。

所有这些交换的合法条件，尽管写法各异，防的是同一件事。

\begin{center}
\begin{tikzpicture}[x=1cm,y=0.7cm,font=\small,>=stealth]
  % ---- 左：尖峰变窄，逐点趋于零，面积恒为 1 ----
  \draw[black!35,->] (-0.2,0) -- (3.9,0);
  \draw[black!35,->] (0,0) -- (0,4.4);
  \draw[blue!45!white,thick] (0,0) -- (1.5,1) -- (3,0);
  \draw[blue!60!black,thick] (0,0) -- (0.75,2) -- (1.5,0);
  \draw[blue!80!black,very thick] (0,0) -- (0.375,4) -- (0.75,0);
  \node[font=\footnotesize,black!70,anchor=west] at (1.0,3.6) {面积恒为 \(1\)};
  \node[font=\footnotesize,black!70,anchor=north,align=center] at (1.8,-0.5)
    {每点趋于 \(0\)，积分不趋于 \(0\)};

  % ---- 右：质量整体溜向尾部 ----
  \begin{scope}[shift={(5.9,0)}]
    \draw[black!35,->] (-0.2,0) -- (5.2,0);
    \draw[black!35,->] (0,0) -- (0,4.4);
    \draw[blue!35!white,thick] (0.2,0) -- (0.8,1.6) -- (1.4,0);
    \draw[blue!60!white,thick] (1.6,0) -- (2.2,1.6) -- (2.8,0);
    \draw[blue!80!black,very thick] (3.0,0) -- (3.6,1.6) -- (4.2,0);
    \draw[black!60,->] (0.9,2.4) -- (3.5,2.4);
    \node[font=\footnotesize,black!70,anchor=south] at (2.2,2.5) {质量整体外移};
    \node[font=\footnotesize,black!70,anchor=north,align=center] at (2.2,-0.5)
      {每点趋于 \(0\)，质量一点没少};
  \end{scope}
\end{tikzpicture}

\emph{两种质量逃逸，一个道理。函数列在每一点上都老老实实趋于零，而积分岿然不动，因为面积要么挤进了越来越窄的尖峰，要么整体溜到了越来越远的地方。极限与积分交换失败，失败的形态就是这两张图。}
\end{center}

\hyperref[sec:12-Real-Analysis-Support/04-Limits-and-Integration/01]{``点态极限为何不能直接穿过积分号''}把这两个反例摆在最前面，理由正是它们把``条件''这件事从形式主义变成了具体的东西：你要防的不是抽象的``不严格''，是这个尖峰和这次外移。\hyperref[sec:12-Real-Analysis-Support/04-Limits-and-Integration/02]{``控制收敛用一个可积包络守住质量''}给出的条件——存在一个不依赖 \(n\) 的可积函数把整列函数压在下面——之所以有效，就是因为一个固定的可积包络同时封住了这两条逃逸路线：尖峰长不了那么高，尾部也放不下那么多面积。

\subsection{同一个条件的其它写法}

一致收敛也在做这件事。\hyperref[sec:12-Real-Analysis-Support/02-Convergence/02]{``点态与一致收敛只差一个量词顺序''}讲的那个量词顺序——每个点各有各的 \(N\)，还是一个 \(N\) 管住所有点——正是尖峰得以存在的空间：点态收敛允许``尖峰所在的那个点''随 \(n\) 移动，一致收敛不允许。紧致性做的是同一件事的拓扑版本，\hyperref[sec:12-Real-Analysis-Support/03-Topology-Basics/02]{``紧致性把无限局部信息压成有限控制''}说明它为什么总在极限论证里出现：无限个局部的界取下确界可能是零，有限个不会；而质量能逃逸，靠的就是``可以一直往外走''这块无限的余地。\hyperref[sec:12-Real-Analysis-Support/05-Measure-Integration-and-Conditional-Expectation/06]{``一致可积性阻止期望质量逃逸''}的名字已经把结论写在标题上了。可积包络、一致有界、紧致——三种语言，一件事。

导数穿过积分号是同一个问题多加一层。\hyperref[sec:12-Real-Analysis-Support/04-Limits-and-Integration/03]{``求导穿过积分号需要控制差商''}指出要控制的不是被积函数本身而是它的差商，因为求导本来就是一个极限。这一条在机器学习里天天被用到：\(\nabla_\theta\E[\ell]\) 与 \(\E[\nabla_\theta \ell]\) 的交换，是随机梯度这整套做法能够成立的前提。\hyperref[sec:14-Deep-Learning/06-Variational-Autoencoders/02]{``重参数化让随机节点可以反向传播''}讲的正是怎样把这个交换变成合法的：把随机性从参数里挪到一个与参数无关的噪声源上，期望的积分区域就不再随参数变动，梯度可以安心穿过去。这不是一个实现技巧，是一次把条件补齐的操作。

累次积分换序有自己的名字。\hyperref[sec:12-Real-Analysis-Support/06-Product-Measures-and-Fubini/02]{``Fubini 与 Tonelli：二重积分何时能拆成两次''}给出两条路：被积函数非负时可以直接换（Tonelli），否则要先验证绝对可积（Fubini）。非负那一条格外好用，因为``质量不能是负的''本身就堵住了正负两部分互相抵消所掩盖的发散。求和换序、期望与求和换序都是它的特例——一个条件收敛的级数重排能收敛到任意指定的数，这个经典事实与积分换序失败是同一个现象。

Jensen 不等式常被当成一条独立的技巧，其实它是这一篇的正面陈述。\hyperref[sec:05-Probability/07-Transforms-and-Bounds/03]{``Jensen 不等式与随机凸性''}说的是期望与凸函数\emph{不能}随便交换，而且给出了不能交换时偏差的方向：\(\E[\varphi(X)]\ge\varphi(\E[X])\)。开头那两版仿真代码就是它的一个实例，取最大是凸的。同一个不等号在强化学习里造成一个具体的、有名字的偏差：\hyperref[sec:16-Control-and-Reinforcement-Learning/02-Reinforcement-Learning/04]{``从动态规划到 Q 学习''}里的最大化偏差——用同一批带噪估计既选出最大的那个动作、又报告它的值，报告出来的期望必然高于真值，因为 \(\E[\max]\ge\max\E\)。它不随样本增多而消失得很快，也不能靠调超参数解决，只能靠把``选''和``评''拆到两批数据上（这正是双重估计的做法）。

随机逼近里还有一处更隐蔽的交换：迭代收敛到的极限，与每步噪声的期望，两者的顺序。每步用的是无偏估计不足以保证极限是无偏的，因为``取极限''与``取期望''之间同样隔着一条缝，需要额外的一致可积或有界性条件把它焊上。\hyperref[sec:13-Machine-Learning/08-Sampling-and-Inference/01]{``Monte Carlo 用随机样本逼近期望''}里那些偶尔出现的巨大样本值，正是这条缝的实际形态——重要性采样的权重方差无界时，样本均值仍然``几乎必然''收敛，但收敛得慢到任何有限次实验都看不出来，而报出的标准误会小得离谱。

最后一处几乎不被当成换序问题：两个极限之间的交换。\hyperref[sec:02-Calculus/06-Multivariable-Calculus/02]{``偏导数为何不等于可微''}里那个沿各方向偏导数都存在却不可微的例子，本质就是``先沿 \(x\) 趋近再沿 \(y\)''与反过来给出不同答案。混合偏导数相等需要连续性条件，也是同一件事。

\subsection{最该记住的那个反例}

上面所有条件都是关于``被积函数别太野''的。有一类失效与被积函数无关，纯粹来自积分区域随参数移动，它在实践中出现得最频繁，而框架永远不会报错。

设 \(X\sim\Uniform(0,\theta)\)，考虑 \(m(\theta)=\E[X^{2}]=\theta^{2}/3\)。直接对 \(\theta\) 求导得 \(2\theta/3\)。若按``把梯度塞进期望里''去算，得到的是 \(\E[\partial_\theta X^{2}]=0\)，因为 \(X^{2}\) 里根本没有 \(\theta\)。两个答案差了整整一项，而那一项来自密度的支撑端点随 \(\theta\) 移动——参数不只影响被积的东西，还影响积分在哪里积。

这个例子值得记住，是因为它的形态很常见：任何``参数决定了取值范围''的模型都有这个问题——截断分布的截断点、生存分析里的删失时间、带硬约束的采样、按阈值决定是否触发的仿真逻辑。写成代码之后，自动微分会尽职地对被积表达式求导，然后返回一个漏掉了边界项的梯度。它不会警告，因为从计算图的角度看一切正常。检查它的办法很朴素：拿数值差分对一两个坐标校验一次。这不是定理，是一条应当写进流程的经验做法。

\subsection{收束}

这一单元把全书横着切了五刀。线性化把弯的暂时当成直的，换基把纠缠的作用变成独立的缩放，迭代把解不出来的解写成一列近似的极限，不变量用一个不动的数一举排除大量可能性，换序在两个操作之间挪动顺序。它们出现在完全不同的学科里，用完全不同的词汇被讲述，但每一个都遵循同一个句式：\emph{先把问题换成一个便宜的形式，再交代这次换算在什么范围内有效。}

这五个动作不构成一份清单，也不该被当成清单来用。它们的价值在识别：读到一段新材料时能认出``这里在做线性化''，接下来该问什么就是确定的了——可信邻域有多宽，我的步子有没有迈出去。认出``这里在迭代''，该问的是在什么范数下压缩、不动点是不是要的解。认出``这里换了个顺序''，该问的是质量会不会跑掉。识别节省的不是记忆量，是在陌生材料面前决定该怀疑什么的时间。

\hyperref[ch:001]{``为什么是 Beyond Formulas''}里说过这本书要处理的困境：每一章读完都懂，合起来不知道自己在干什么。走到这里可以把那句话补完。这本书从头到尾给的不是方法——方法在每一本教材里都有，而且写得比这里详细。它给的是判断方法何时可信的依据：这个定理的哪条假设挡住了哪种失效，这个近似在什么尺度上还算数，这个保证承诺到哪一层为止。

而这些依据，追到底并没有多少条。它们反复归结为同一个问题的不同问法——\emph{这个动作，在什么条件下合法。}前面十六个部分是这个问题在十六个学科里的具体答案，这一单元是把答案叠在一起之后剩下的那个形状。下次遇到一个没学过的方法、一篇看不懂的论文、一个跑出奇怪数字的模型，能派上用场的多半不是记住的公式，是这个问题本身。
